\section{Actuation and propulsion}
\label{sec:actuation-propulsion}

The canonical multirotor plant integrates actuator dynamics and propulsion dynamics inside one RHS (\code{Sim.PlantModels.CoupledMultirotorModel}) so that both fixed-step and adaptive integrators see a single continuous system.

\subsection{Actuator channel mapping and dynamics}

PX4 produces actuator commands as fixed-size ABI arrays: 12 motor channels in $[0,1]$ and 8 servo channels in $[-1,1]$ (\code{Vehicles.ActuatorCommand}). These are \emph{normalized lockstep ABI units} (not PWM) taken from the uORB outputs subscribed by the bridge (\code{actuator\_motors} and \code{actuator\_servos}); they are interpreted as normalized ESC duty and control-surface deflection. The plant stores actuator outputs in the same fixed-size arrays ($\vect{y}_m\in\mathbb{R}^{12}$, $\vect{y}_s\in\mathbb{R}^{8}$), but the physical vehicle may have $N\le 12$ propulsors.

\paragraph{Mapping from PX4 channels to physical motors.}
A typed mapping (\code{Vehicles.MotorMap\{N\}}) selects which PX4 motor channel drives each physical motor:
\begin{equation}
  d_i \coloneqq (\vect{y}_m)[\texttt{motor\_channel}[i]],\qquad i\in\{1,\dots,N\}.
\end{equation}
Faults are defined over \emph{physical} motors (indices $1..N$), not PX4 channels. The plant applies the motor-disable mask by forcing the affected physical duties to zero before evaluating propulsion.

\paragraph{Actuator dynamics.}
Actuator channel dynamics are integrated as part of the plant state. Motor and servo channels can each use one of the following models (\code{src/sim/Vehicles.jl}, \code{src/sim/PlantModels/CoupledMultirotor.jl}):
\begin{itemize}[leftmargin=*,itemsep=2pt]
  \item \textbf{Direct}: algebraic output, snapped to the new command at each autopilot boundary via \code{plant\_on\_autopilot\_tick}, then held.
  \item \textbf{First-order}: \(\dot{\vect{y}} = (\vect{u}_{cmd} - \vect{y})/\tau\).
  \item \textbf{Second-order}: \(\dot{\vect{y}}=\dot{\vect{y}}_{\mathrm{eff}}\), \(\ddot{\vect{y}}=\omega_n^2(\vect{u}_{cmd}-\vect{y})-2\zeta\omega_n\dot{\vect{y}}_{\mathrm{eff}}\), with an optional rate limit implemented by clamping \(\dot{\vect{y}}_{\mathrm{eff}}\) in the RHS and projecting stored rate states post-step.
\end{itemize}

After each engine interval the coupled plant implements \code{plant\_project} to clamp actuator outputs into ABI ranges and to project second-order rate states back into the configured rate limit, ensuring hard bounds are enforced deterministically.

\begin{figure}[t]
  \centering
  \includegraphics[width=\linewidth]{figs/actuator_step_response.png}
  \caption{Actuator step response example: direct vs second‑order dynamics with a rate limit.}
  \label{fig:actuator-step-response}
\end{figure}

\begin{figure}[t]
  \centering
  \includegraphics[width=\linewidth]{figs/motor_spinup.png}
  \caption{Motor/prop spin‑up example for a step duty command at constant bus voltage.}
  \label{fig:motor-spinup}
\end{figure}

\subsection{Motor/ESC/prop model}
\label{sec:propulsion-model}

Each physical motor uses a parameter-only container \code{Propulsion.MotorPropUnit} consisting of:
(i) an ESC stage with efficiency $\eta$ and deadzone, (ii) a quasi-static BLDC electrical model, and (iii) a quadratic propeller model with an inflow correction. The plant integrates rotor speed $\Omega_i$ (rad/s) for each physical motor.

\paragraph{ESC stage.}
The commanded duty is clamped to $[0,1]$ and zeroed if it is below the ESC deadzone. The motor terminal voltage is
\begin{equation}
  V_m = d\,V_{bus}.
\end{equation}

\paragraph{BLDC electrical model (quasi-static current).}
Let $K_v$ be specified in RPM/V. The back-EMF constant is
\begin{equation}
  K_e = \frac{60}{2\pi K_v} \quad \left[\mathrm{V}\,\mathrm{s}/\mathrm{rad}\right],
\end{equation}
and the torque constant is set equal to $K_t\equiv K_e$ (SI). Motor current is computed from applied voltage and back-EMF and then clamped:
\begin{equation}
  I_m = \mathrm{clamp}\!\left(\frac{V_m - K_e\,\Omega}{R},\; 0,\; I_{\max}\right).
\end{equation}
The electrical torque subtracts an approximate no-load current $I_0$:
\begin{equation}
  \tau_e = K_t\,\max(0, I_m - I_0).
\end{equation}

\paragraph{Loss modeling note.}
In this model, $I_0$ represents \emph{electrical} no-load current (iron/core loss) that does not produce torque, while $b\,\Omega$ is a \emph{mechanical} drag torque. The defaults are a pragmatic generic fit rather than a motor identification; if you want a single loss channel, set either $I_0=0$ (pure mechanical drag) or $b=0$ (pure electrical loss) to avoid double-counting.

\paragraph{Quadratic propeller aerodynamics.}
At air density $\rho$ the thrust and load torque are
\begin{equation}
  T = k_T\,\rho\,\Omega^2\,f_T(\mu),\qquad
  Q = k_Q\,\rho\,\Omega^2\,f_Q(\mu),
\end{equation}
with a simple inflow sensitivity factor
\begin{equation}
  f(\mu) = \frac{1}{1 + k\,\mu^2},\qquad \mu = \frac{V_{ax}}{\max(10^{-3}, |\Omega|R_{prop})}.
\end{equation}
The prop radius $R_{prop}$ is stored in the prop parameters. The coupled plant computes the axial component $V_{ax}$ by projecting the \emph{air velocity relative to the vehicle} onto each propulsor axis (\cref{sec:axis-convention}):
\begin{equation}
  \vect{v}_{air,n} = \vect{w}_{ned} - \vect{v}_{ned},\qquad
  \vect{v}_{air,b} = \Rbn^{\mathsf{T}}\,\vect{v}_{air,n},\qquad
  V_{ax,i} = -\vect{v}_{air,b}\cdot \vect{a}_{b,i}.
\end{equation}
The inflow correction uses $\mu^2$, so the sign of $V_{ax}$ affects neither $f_T$ nor $f_Q$. This $1/(1+k\mu^2)$ attenuation is an ad-hoc shaping term (not a blade-element or momentum-theory model) chosen for stability and simplicity; it cannot represent windmilling, vortex-ring, or braking regimes. See \cref{sec:design-choices}.

\paragraph{Rotor-speed dynamics.}
In the coupled plant, rotor speed is a continuous state integrated by the same ODE integrator as the rest of the plant. The RHS evaluates
\begin{equation}
  \dot{\Omega} = \frac{\tau_e - Q - b\,\Omega}{J},
\end{equation}
where $J$ is the axial inertia and $b$ is a viscous friction coefficient, and where $Q$ is the aerodynamic load torque evaluated at the \emph{current} $\Omega$ (no internal Euler step inside the RHS). A guard prevents unphysical negative spin: if $\Omega\le 0$ and $\dot{\Omega}<0$, the derivative is set to zero.

\paragraph{Bus current and signed reaction torque.}
The current drawn from the bus is
\begin{equation}
  I_{bus} = \frac{d\,I_m}{\max(10^{-6},\eta)}.
\end{equation}
The quadratic prop model returns a positive load torque magnitude $Q\ge 0$. The coupled plant applies a per-rotor sign \code{rotor\_dir[i]} to generate the signed reaction torque on the body,
\begin{equation}
  Q_{body,i} = \texttt{rotor\_dir}[i]\,Q_i,
\end{equation}
and passes this signed value into the rigid-body dynamics.

\subsection{Default propulsion parameterization and calibration}
\label{sec:propulsion-defaults}

The helper \code{Propulsion.default\_multirotor\_set} provides a generic baseline motor+prop set. It calibrates the thrust coefficient $k_T$ by solving for the steady-state speed at duty $d=1$ under a quadratic load and then bisection-searching $k_T$ so that the resulting thrust matches a target.

In the aircraft builder (\code{src/sim/Aircraft/Build.jl}), this target is derived from the airframe mass and gravity:
\begin{equation}
  T_{hover,rotor} = \frac{m\,g}{N},\qquad T_{target} = \texttt{thrust\_calibration\_mult}\cdot T_{hover,rotor},
\end{equation}
with \texttt{thrust\_calibration\_mult}=2.0 by default. The drag torque coefficient is tied to $k_T$ through \texttt{km\_m} via $k_Q = \texttt{km\_m}\,k_T$.

\paragraph{Caveat.}
This calibration is a pragmatic way to get reasonable thrust scaling for a ``generic'' vehicle. It is not an identification procedure for a specific airframe; if you care about realism beyond closed-loop controllability, replace the parameters with airframe-specific values.


\subsection{TOML configuration: actuation mapping and propulsion parameters}
\label{sec:actuation-toml}

Actuation and propulsion are configured in two places in the aircraft spec:
\begin{itemize}[leftmargin=*,itemsep=2pt]
  \item \code{[actuation]} defines how PX4 ABI channels map to physical motors/servos and which actuator dynamics model is used.
  \item \code{[airframe.propulsion]} defines the motor/ESC/prop parameters used by the coupled plant (\cref{sec:propulsion-model}).
\end{itemize}

A minimal Iris-like configuration is:

\begin{lstlisting}[language=toml]
[actuation]

[[actuation.motors]]
id = "motor1"
channel = 1

[[actuation.motors]]
id = "motor2"
channel = 2

[[actuation.motors]]
id = "motor3"
channel = 3

[[actuation.motors]]
id = "motor4"
channel = 4

[actuation.motor_actuators]
kind = "direct"   # direct | first_order | second_order

[actuation.servo_actuators]
kind = "direct"

[airframe.propulsion]
kind = "multirotor_default"
km_m = 0.05
V_nom = 12.0
rho_nom = 1.225
rotor_radius_m = 0.127
inflow_kT = 8.0
inflow_kQ = 8.0
thrust_calibration_mult = 2.0
rotor_dir = [1, 1, -1, -1]

[airframe.propulsion.esc]
eta = 0.98
deadzone = 0.02

[airframe.propulsion.motor]
kv_rpm_per_volt = 920.0
r_ohm = 0.25
j_kgm2 = 1.2e-5
i0_a = 0.6
viscous_friction_nm_per_rad_s = 2.0e-6
max_current_a = 60.0
\end{lstlisting}

\paragraph{Actuation mapping semantics.}
Each entry in \code{[[actuation.motors]]} creates a named physical motor ID and assigns it a PX4 channel index. During build, IDs are resolved into a \code{MotorMap\{N\}} so that \(\texttt{channel}\) implements the lookup in \cref{sec:actuation-propulsion}:
\[
  d_i \leftarrow (\vect{y}_m)[\texttt{channel}_i].
\]
Unknown IDs referenced elsewhere (e.g. by buses) are rejected at build time.

\paragraph{Actuator model parameters.}
The actuator model may be given as a string shorthand (\code{kind="direct"} etc) or as a table with parameters. For example, a second-order motor actuator with a finite rate limit uses:
\begin{lstlisting}[language=toml]
[actuation.motor_actuators]
kind = "second_order"
"\u03C9n" = 20.0   # omega_n
"\u03B6"  = 0.7    # zeta
rate_limit = 100.0   # units: (duty)/s
\end{lstlisting}
These keys map directly to \(\omega_n\), \(\zeta\), and the optional \(\dot{\vect{y}}\) clamp described in \cref{sec:actuation-propulsion}.

\paragraph{Propulsion parameter mapping.}
The propulsion keys match the symbols used in \cref{sec:propulsion-model}:
\begin{itemize}[leftmargin=*,itemsep=2pt]
  \item \code{kv\_rpm\_per\_volt} \(\rightarrow K_v\) (and thus \(K_e=60/(2\pi K_v)\) and \(K_t=K_e\)).
  \item \code{r\_ohm} \(\rightarrow R\) in \(I_m=(V_m-K_e\Omega)/R\).
  \item \code{i0\_a} \(\rightarrow I_0\) in \(\tau_e=K_t\max(0,I_m-I_0)\).
  \item \code{max\_current\_a} \(\rightarrow I_{\max}\) current clamp.
  \item \code{j\_kgm2} and \code{viscous\_friction\_nm\_per\_rad\_s} \(\rightarrow J\) and \(b\) in \(\dot{\Omega}=(\tau_e-Q-b\Omega)/J\).
  \item \code{eta} and \code{deadzone} configure the ESC stage (\cref{sec:propulsion-model}).
  \item \code{rotor\_radius\_m} \(\rightarrow R_{prop}\) in the inflow ratio \(\mu\).
  \item \code{inflow\_kT}, \code{inflow\_kQ} \(\rightarrow k\) in \(f(\mu)=1/(1+k\mu^2)\) for thrust and torque respectively.
  \item \code{km\_m} ties \(k_Q\) to \(k_T\) through \(k_Q=\texttt{km\_m}\,k_T\) in the default prop model.
  \item \code{thrust\_calibration\_mult} sets the target thrust multiplier used during build-time calibration (\cref{sec:propulsion-defaults}).
  \item \code{rotor\_dir[i]} signs the reaction torque applied to the body, \(Q_{body,i}=\texttt{rotor\_dir}[i]\,Q_i\). Rotor spin sign is the opposite (\(s_i=-\texttt{rotor\_dir}[i]\)) for gyroscopic coupling in \cref{sec:rigid-body}.
\end{itemize}
